\documentclass[10pt,a4paper]{amsart}
\usepackage[margin=19mm]{geometry}
\usepackage{amsmath,amssymb,mathtools}
\usepackage[scr=rsfs,cal=euler]{mathalfa}
\usepackage{xfrac}
\usepackage[unicode,hidelinks,bookmarks=false]{hyperref}
\newcommand{\ie}{i.e.\ }
\newcommand{\cf}{cf.\ }
\newcommand{\resp}{respectively\ }
\newcommand{\Subsection}{\subsection}
\providecommand{\nobreakdash}{\nobreak\hbox{-}}
\setlength{\emergencystretch}{2em}
\multlinegap0pt
\usepackage{amssymb}
\newtheorem{theorem}{Theorem}
\newtheorem{proposition}{Proposition}
\newtheorem{lemma}{Lemma}
\def \R{{\mathbb R}}
\def \s{{\mathbb S}}
\title{Existence and stability of weak critical points of $r$-energy functionals}
\author{Stefano Montaldo, Andrea Ratto, Antonio Sanna}
\date{Source: arXiv:2605.03532v1 (CC BY 4.0)}
\begin{document}
\maketitle

\noindent\emph{Source and licence.} Existence and stability of weak critical points of $r$-energy functionals. Version arXiv:2605.03532v1 (CC BY 4.0), distributed by arXiv under the Creative Commons Attribution 4.0 International licence, http://creativecommons.org/licenses/by/4.0/, as recorded in the arXiv licence metadata for this exact version and shown on the abstract page https://arxiv.org/abs/2605.03532v1. Full licence notice: \texttt{licenses/arxiv-2605.03532-CC-BY-4.0.txt}.\par\smallskip

\noindent\textit{Editorial scope.} Counted unit: the article's theorem th-ellipsoid-triharm (source0.tex lines 230--235). The article's complete original proof of it is delivered with it (source0.tex lines 792--861, one \texttt{proof} environment of 70 lines, which the article places away from the statement). No mathematics is written, repaired or re-derived: the statement and the proof are the article's own lines, in the article's own notation, and no retained line is substituted or re-flowed.  Retained uncounted as the source-local context the proof needs: lines 151--156; lines 212--217; lines 312--314; lines 318--321; lines 412--414; lines 751--756; lines 764--766.  Nothing was omitted from any retained span, so the package carries no omission record.  No retained line contains a citation, so the wrapper carries no bibliography and no bibliography entry.  No retained line contains a figure environment, an image-inclusion command or drawing code, so no image is delivered and the package declares no asset dependency.  Editorial formatting only, in addition: an AMS \texttt{amsart} preamble that restates the article's own declarations and macros used by the retained lines, namely the environments \texttt{theorem}, \texttt{proposition}, \texttt{lemma} on separate counters, together with the standard packages those lines need.  The retained environments print in this document as Theorem 1, Proposition 1, Lemma 1 in order of appearance, whereas the article prints the same items with the numbering of its own full document.  The complete native source of the article as distributed in the arXiv e-print for this exact version is bundled unchanged as \texttt{sources/199794/source0.tex}.
\begin{equation}\label{symmetricmaps}\begin{array}{lcll}
                                           \varphi_{a} \,:\, &B^n &\to &  \s^n \subset \R^n \times \R \\
                                           &&& \\
                                           &x &\mapsto & ( \sin a\,{\frac{x}{\rho}}, \,\cos a) \, ,
                                         \end{array}
\end{equation}

\begin{equation}\label{symmetricmaps-ellipsoid}\begin{array}{lcll}
                                           \varphi_{a} \,:\, &B^n &\to &  E^n(b) \subset \R^n \times \R \\
                                           &&& \\
                                           &x &\mapsto & ( \sin a\,{\frac{x}{\rho}}, b\,\cos a) \, ,
                                         \end{array}
\end{equation}

\begin{proposition}\label{pro:weakly-criterium} Let
$\varphi_a:B^n \to \s^n$ be a map as in \eqref{symmetricmaps}. Then $\varphi_a$ is a weak critical point of $E^{ES}_r(\varphi)$ ($E_r(\varphi)$) if and only if  $\varphi_a \in W^{r,2}\left ( B^n,\s^n \right )$ and $\varphi_a$ is $ES-r$-harmonic ($r$-harmonic) on $B^n \setminus \{ O\}$.
\end{proposition}

\begin{lemma}\label{lemma-belong-Sobolev-space}
Let
$\varphi_a:B^n \to \s^n$ be a map as in \eqref{symmetricmaps}. Then $\varphi_a \in W^{r,2}\left ( B^n,\s^n \right )$ if and only if $n\geq 2r+1$.
\end{lemma}

\begin{equation}\label{Euler-lagrange-generale}
\sum_{i=1}^r \, (-1)^i \, \frac{d^i}{d\rho^i}\,\left( \frac{\partial L_r}{\partial \alpha^{(i)}}\right)+\frac{\partial L_r}{\partial \alpha}=0 \,\,.
\end{equation}

\begin{equation}\label{symmetricmaps-ellipsoid-general}\begin{array}{lcll}
                                           \varphi_{\alpha} \,:\, &B^n &\to &  E^n(b) \subset \R^n \times \R \\
                                           &&& \\
                                           &x &\mapsto & ( \sin \alpha(\rho)\,\frac{x}{\rho}, b\,\cos \alpha(\rho)) \, .
                                         \end{array}
\end{equation}

\begin{equation}\label{eq-tau-Enb}
\tau_\alpha= \ddot{\alpha}+\frac{(n-1)}{\rho}\dot \alpha-\frac{(n-1)}{\rho^2} \frac{\sin \alpha \cos \alpha}{k^2(\alpha)}+\frac{k'(\alpha)}{k(\alpha)} \dot{\alpha}^2 \,,
\end{equation}

\begin{theorem}\label{th-ellipsoid-triharm}
Assume  $n\geq 7$.  Then there exists a proper weakly triharmonic map $\varphi_a:B^n \to E^n(b)$ of type \eqref{symmetricmaps-ellipsoid} if and only if
\begin{equation}\label{eq-cond-trihar-ellipsoid}
0<b^2<\frac{(n-1)}{4(n-6)}\,.
\end{equation} 
\end{theorem}

\begin{proof}[Proof of Theorem~\ref{th-ellipsoid-triharm}]
First, using the argument of Lemma~\ref{lemma-belong-Sobolev-space}, we observe that $\varphi_a \in  W^{3,2}\left ( B^n,E^n(b) \right )$ if and only if $n \geq 7$.
Also in this case, we observe that a map $\varphi_a:B^n \to E^n(b)$ as in \eqref{symmetricmaps-ellipsoid} belongs to the family of rotationally symmetric maps described by \eqref{symmetricmaps-ellipsoid-general}. Moreover, by the argument of Proposition~\ref{pro:weakly-criterium}, it is sufficient to check when $\varphi_{\alpha}$ restricted to $B^n\setminus\{O\}$ is triharmonic. Now a routine computation shows that the trienergy of these maps is given, up to an irrelevant constant factor, by
\begin{equation}\label{eq-bienergy-ellipsoid}
E_3\left( \varphi_\alpha \right )=\int_0^1 L_3\,d\rho \,,
\end{equation}
where
\[
L_3= \rho^{n-1} \left((n-1) \frac{\tau^2_{\alpha}  \cos^2 \alpha}{\rho^2}+ {\dot\alpha}^2\left(\tau_{\alpha}k'(\alpha)+k(\alpha) \frac{\partial}{\partial\alpha}{{\tau}_{\alpha}}\right)^2\right)\,
\]
where $\tau_\alpha$ and $k(\alpha)$ where defined in \eqref{eq-tau-Enb}. Now, using \eqref{Euler-lagrange-generale}, one can formally compute the triharmonicity equation. Although the full expression of the triharmonicity ODE is rather complicated, in the special case that $\alpha(\rho)\equiv a$ the triharmonicity condition, in $x=\cos^2 a$, reduces to the polynomial equation
\begin{equation}\label{eq-triharmonicity-condition-ellipsoid}
P_{b}(x)=A_3 x^3+A_2 x^2+A_1 x+A_0=0 \,,
\end{equation}
where
{\small
\[
\begin{split}
A_3=&(b^2-1)  \left(2(n-4) b^2-3 (n-3)\right) \left(4 (n-6) b^2-5 (n-5)\right)\\
A_2=&-b^2 \left(24 (n-4) (n-6)b^4 + (-62 n^2+566 n-1224)b^2+41 n^2-332 n+651\right)\\
A_1=&2 b^2 \left(12 (n-4) (n-6) b^4+(-17 n^2+149 n-312)b^2+(n-1) (2 n-5)\right)\\
A_0=&-2 b^4 (n-4) \left(4(n-6) b^2-n+1\right)\,.
\end{split}
\]
}
Thus the conclusion of the theorem is equivalent to the statement that the third order polynomial $P_b(x)$ defined in \eqref{eq-triharmonicity-condition-ellipsoid} admits a root in the interval $(0,1)$  if and only if \eqref{eq-cond-trihar-ellipsoid} holds. Now, since $n \geq 7$, $P_b(1)=-15 \left(n^2-8 n+15\right)<0$. Moreover, if \eqref{eq-cond-trihar-ellipsoid} holds, then $P_b(0)>0$ and so there exists a root of $P_b(x)$
in the interval $(0,1)$.

Conversely, let us assume that \eqref{eq-cond-trihar-ellipsoid} does not hold, that is
\[
b^2\geq\frac{(n-1)}{4(n-6)}=b_*^2\,.
\]
We prove that in this case $P_{b}(x)$ is strictly negative for $x \in(0,1)$. To this purpose, first we show that $P_{b_*}(x)<0$ for all $x\in (0,1)$ and $n\geq 7$. Indeed, we have 
\begin{eqnarray*}
P_{b_*}(x)&=& \frac{x}{4(n-6)}\Big [-3 (n-5) (n-1)^2-(n-1) \left(27 n^2-268 n+601\right) x\\
&&-2 (3 n-23) \left(5 n^2-49 n+104\right) x^2 \Big ]\,.
\end{eqnarray*}
Now, direct inspection shows that $P_{b_*}(x)<0$ for $x\in (0,1)$ and $n\geq 8$. If $n=7$, $P_{b_*}(x)$ becomes
\[
\frac{x}{4}(24 x^2-288 x-216)
\]
which is negative for all $x\in (0,1)$. Now, in summary, to end the proof it is sufficient to show that for $b^2> b_*^2$ we have
\[
P_{b_*}(x)-P_{b}(x)\geq 0\,.
\]
A direct computation yields
\[
P_{b_*}(x)-P_{b}(x)=\frac{(1-x) (4 (n-6)b^2+(1-n))}{4 (n-6)} Q_b(x)\,,
\]
where
\begin{eqnarray*}
Q_b(x)&=&x \left(30 n^2 x+3 n^2-286 n x-18 n+616 x+15\right)\\
&&+4 (n-6) (7 n-25) (1-x)\, x\, b^2+8 (n-6) (n-4) (1-x)^2 b^4\,.
\end{eqnarray*}

Thus it remains to show that $Q_b(x)\geq 0$ for $x\in (0,1)$, $n\geq 7$ and $b^2> b_*^2$.

Indeed, since $n\geq 7$, $b^2> b_*^2$ implies $Q_b(x)>Q_{b_*}(x)$, where
\[
\begin{split}
Q_{b_*}(x)=&\frac{1}{2(n-6)}\Big[((n-4) (n-1)^2+2 (n-4) (n-1) (9 n-59) x\\
&+\left(47 n^3-790 n^2+4239 n-7096\right) x^2\Big]\,.
\end{split}
\]
Now, direct inspection shows that $Q_{b_*}(x)>0$ when $n\geq 8$ and $x\in (0,1)$. Finally, if $n=7$, then $Q_{b_*}(x)$ becomes
\[
\frac{-12 x^2+144 x+108}{4}
\]
which is clearly positive for all $x\in (0,1)$ and so the proof is completed.
\end{proof}

\end{document}
